% =====================================================================
% 07 — Inference methods for EBMs.
% =====================================================================
% Planning (% comments only):
% Langevin / SGLD, HMC, MPPI, Gibbs (discrete), MAP-via-gradient,
% energy descent at inference (adaptive compute), variational inference,
% diffusion ancestral sampling, MCMC sampling for ratios.
% Endpoint: adaptive-compute inference is the load-bearing claim that
% EBMs add beyond AR LLMs — but only if iteration buys more than equivalent
% additional training.
% =====================================================================

\section{Inference in Energy-Based Models}
\label{sec:ebm_inference}

Sampling from \cref{eq:ebm_definition} or computing posterior queries against it is the second place the partition function bites.
This section organizes inference methods by the discreteness of the object and by whether the model exposes its score, its energy, or only its conditional updates.

\subsection{Continuous: Langevin dynamics and Hamiltonian methods}
\label{ssec:langevin}

Langevin dynamics defines an SDE whose stationary distribution is $\ptheta$:
\begin{equation}
  \mathrm{d} x \;=\; -\nabla_x \Etheta(x)\, \mathrm{d}t \;+\; \sqrt{2}\, \mathrm{d}W_t,
  \label{eq:langevin}
\end{equation}
discretized as $x_{k+1} = x_k - \tau \nabla_x \Etheta(x_k) + \sqrt{2\tau}\, \eta_k$ for noise $\eta_k \sim \mathcal{N}(0, I)$ and step $\tau > 0$.
Convergence to $\ptheta$ holds under mild regularity~\citep{roberts1996exponential}.
The cost of inference is the number of iterations to traverse a representative region of the energy landscape.
Stochastic gradient Langevin dynamics~\citep{welling2011bayesian} couples \cref{eq:langevin} with stochastic gradient updates to scale.
Hamiltonian Monte Carlo~\citep{neal2011mcmc,betancourt2017conceptual} augments with momentum, reducing the number of effective steps when the energy is smooth.

Score-based generative models~\citep{song2019generative,song2021scorebased} use a noise-conditional score network in place of $\nabla_x \Etheta$ and integrate \cref{eq:langevin} along a denoising schedule.
The training methods of \cref{ssec:score_matching} return an estimate of the score that this inference loop consumes.

\subsection{Discrete: Gibbs, simulated annealing, and local-search analogues}
\label{ssec:gibbs}

For discrete sequence spaces, Langevin dynamics does not directly apply.
Gibbs sampling~\citep{geman1984stochastic} updates one coordinate at a time conditional on the rest;
the conditional itself requires evaluating $|\V|$ energies per step.
Block-Gibbs and Metropolis--within-Gibbs variants trade per-step cost for mixing rate.
For energies parameterized as deep networks over discrete sequences, recent work~\citep{grathwohl2021ono,emami2024discrete} proposes continuous relaxations and gradient-guided proposals that accelerate mixing without losing the discrete target.

Simulated annealing~\citep{kirkpatrick1983optimization} traces the same energy under a temperature schedule that contracts as iteration proceeds.
It is a natural model for ``inference as optimization'' when only the MAP is required.

\subsection{Variational inference and the free-energy view}
\label{ssec:variational}

Variational methods replace exact inference with a tractable surrogate $q_\phi(z)$ over latents $z$ and minimize the variational free energy
\begin{equation}
  \mathcal{F}[q_\phi] \;=\; \EE_{q_\phi}\!\left[\Etheta(z)\right] \;-\; H(q_\phi),
  \label{eq:variational_free_energy}
\end{equation}
where $H(q_\phi)$ is the entropy.
Minimization of \cref{eq:variational_free_energy} gives the closest $q_\phi$ to $\ptheta$ in reverse-KL~\citep{wainwright2008graphical}.
\citet{rezende2014stochastic,kingma2014auto} introduce amortized variants that learn $q_\phi$ as a neural network and have become the standard scaffold for variational autoencoders.
\citet{ranganath2014black} make the approach black-box-applicable.
The mean-field assumption that $q_\phi$ factorizes across coordinates is the analogue of independent token modeling and is what attention's row-stochastic structure already implements~\citep{ramsauer2020hopfield}.

\subsection{Posterior queries: conditional sampling and amortization}
\label{ssec:conditional}

Many practical EBM uses do not need samples from the unconditional $\ptheta$.
They need samples from $\ptheta(x \mid c)$ for a conditioning context $c$ --- a prompt, a partial program, a goal state.
\citet{du2023reduce} demonstrate that conditional energy descent is competitive with autoregressive sampling on structured outputs, and that energy-based composition can be applied in latent space.
For tasks with a hard verifier, conditional sampling on the verifier's accept set reduces to constrained energy descent: project at each step onto the verifier-bounded subspace~\citep{cobbe2021training}.

\subsection{Energy descent at inference and adaptive compute}
\label{ssec:adaptive_compute}

The reason inference deserves a dedicated section is the claim that motivated the field in the first place: \emph{the number of inference steps is a free parameter}, and energy-based inference can spend more on harder inputs.
This is the architectural counterpart to chain-of-thought.
Where chain-of-thought converts more inference into more tokens, energy descent converts more inference into more iterations of the same forward pass~\citep{lecun2022autonomous,du2023reduce}.

The empirical question is whether $K > 1$ steps of energy descent on a fixed model are ever more compute-efficient than the equivalent additional training compute spent on a one-step model.
For well-separated patterns under \cref{eq:hopfield_energy}, \citet{ramsauer2020hopfield} prove that one step is already optimal.
For not-well-separated inputs --- precisely the inputs we care about --- multiple steps are expected to help, and partial results in this direction exist~\citep{bachlechner2021rezero,bansal2022end,schwarzschild2021can}.
The companion suite (\texttt{experiments/02\_ar\_pros\_cons/probes/probe\_05.py}) explicitly tests whether iteration is compute-cheaper than parameter count in a verifier-bounded setting.

\subsection{What the inference survey leaves}
\label{ssec:inference_leaves}

Inference in energy-based models is a budgeted decision: how many iterations, with what proposal, against which conditioning set, with what verifier hook.
The autoregressive recipe makes the equivalent decision implicitly --- one forward pass, beam search, or sample-then-verify --- but does so with a fixed compute budget per output token.
The interesting comparison is not ``EBM inference is hard'' versus ``AR inference is easy.''
It is ``where do you want to spend the variable-compute budget,'' phrased clearly enough that the trade is visible.
\Cref{sec:other_alts} broadens the comparison to alternatives that occupy intermediate positions on this trade.
